\problemname{Sågverket}

Alma har en massa brädor som hon vill såga i bitar. Hon har därför gått till ett toppmodernt
sågverk som automatiskt sågar itu och sorterar brädor. Varje bräda kan representeras av en oändligt
lång $x$-axel, och sågverket kommer såga av brädan vid de $N+1$ positionerna $x_1, x_2, \dots, x_N, y$,
där $y$ är ett tal som användaren får välja, medan talen $x_i$ alltid är samma. Därefter sorteras alla
ändligt långa bitar och matas ut av sågverket. 

Alma vill ha reda på talen $x_1, x_2, \dots, x_N$, men det verkar som att ingen vet hur sågverket fungerar.
Så istället tänker hon lista ut talen genom att mata in några brädor med väl valda värden $y$.

Det finns $N$ hemliga heltal $x_1 < x_2 < \dots < x_N$ mellan $1$ och $10^9$. 
Notera att alla dessa tal är olika. Ditt mål är att hitta talen. Du får ställa frågor till sågverket. 
Sågverket tar ett heltal $y$ som input ($1 \leq y \leq 10^9$), och gör följande:

\begin{enumerate}
  \item Skapa listan $L = [x_1, x_2, \dots, x_N, y]$.
  \item Sortera $L$.
  \item Skapa listan $D$, där $D_i = L_{i+1}-L_i$ för alla $i = 1, 2, \dots, N$.
  \item Sortera $D$, och returnera de $N$ talen i $D$.
\end{enumerate}

Du får ställa högst $N$ frågor, förutom i testgrupp $4$ där du får ställa $N+1$ frågor.

\section*{Interaktion}
Ditt program ska först läsa in två heltal $N$ och $T$ ($1 \leq N \leq 1000$, $1 \leq T \leq 5$).
$N$ är antalet hemliga tal du ska hitta, och $T$ är numret på testgruppen. Anledningen till att $T$ 
är givet i input är för att göra det lättare att ta delpoäng.

Därefter får du börja ställa frågor. Skriv ut en rad med ``\verb|? y|'' för att stoppa in talet
$y$ i sågverket. Talet $y$ måste uppfylla $1 \leq y \leq 10^9$. Efter varje fråga ska
ditt program läsa in $N$ heltal på en rad, talen $D_1, D_2, \dots, D_N$. Notera att $D_1$ kan vara
noll, ifall $y$ var lika med ett av talen $x_i$.

När du har hittat talen $x_1, x_2, \dots, x_N$ ska du skriva ut en rad på formen 
$$! \; x_1 \; x_2 \; x_3 \; \dots \; x_N$$
Denna rad räknas inte som en fråga till sågverket. Därefter ska ditt program avsluta och 
inte skriva ut något mer.

\textbf{Se till att flusha outputen efter varje query}, annars kan du få \textit{Time Limit Exceeded}.
I C++ kan detta göras med exempelvis \texttt{cout << flush}
eller \texttt{fflush(stdout)};
i Python med \texttt{stdout.flush()};
och i Java med \texttt{System.out.flush()}.

\section*{Poängsättning}
Din lösning kommer att testas på en mängd testfallsgrupper.
För att få poäng för en grupp så måste du klara alla testfall i gruppen.

\noindent
\begin{tabular}{| l | l | p{12cm} |}
  \hline
  Grupp & Poängvärde & Gränser \\ \hline
  $1$   & $15$       & $N = 1$\\ \hline
  $2$   & $15$       & $N = 2$  \\ \hline
  $3$   & $11$       & $x_i \leq N+1$ \\ \hline
  $4$   & $37$       & $N \leq 100$, $x_i \leq 10^4$, du får ställa högst $N+1$ frågor. \\ \hline
  $5$   & $22$       & Inga ytterligare begränsningar \\ \hline
\end{tabular}